\documentclass[12pt]{article}
\usepackage{amsmath, amssymb, booktabs, geometry}
\geometry{margin=1in}
\setlength{\parskip}{1em}
\setlength{\parindent}{0pt}

\title{Solutions -- Practice Problem Set 3}
\author{ECON 101 -- Macroeconomics}
\date{Fall 2025}

\begin{document}
\maketitle


\section*{Part A – Conceptual Questions}

\textbf{1. Economic Growth vs. Business-Cycle Expansion}

\textbf{Question:}  
Explain the difference between long-run economic growth and short-run business-cycle expansion. Why is sustained growth primarily determined by productivity rather than labour hours?

\textbf{Solution:}  
\begin{itemize}
\item \textbf{Economic growth} refers to the long-term increase in real GDP per capita — that is, sustained increases in the productive capacity of the economy.
\item \textbf{Business-cycle expansion} refers to short-term increases in real GDP as the economy recovers from a downturn or moves toward full employment.
\item Sustained growth depends on productivity (output per worker) because there is a limit to how much total labour hours can increase — workers can only work so many hours. Productivity gains, driven by technological progress and capital accumulation, allow output to grow even when labour hours remain constant.
\end{itemize}



\textbf{2. Productivity and Institutions}

\textbf{Question:}  
List three key factors that determine labour productivity and discuss how the financial system contributes to each.
\newpage
\textbf{Solution:}  
\begin{enumerate}
    \item \textbf{Physical capital:} Machinery, equipment, and buildings that enhance worker efficiency. Financial markets supply funds for firms to invest in new capital.
    \item \textbf{Human capital:} Skills and education that improve worker quality. Financial systems enable student loans and corporate training investments.
    \item \textbf{Technological progress:} Innovation and improved production methods. Banks, stock markets, and venture capital provide financing for R\&D and innovation.
\end{enumerate}
The financial system channels savings into productive investment, ensuring resources are used efficiently to enhance productivity.



\textbf{3. Functions of the Financial System}

\textbf{Question:}  
Describe three key roles of the financial system in promoting economic growth.

\textbf{Solution:}  
\begin{itemize}
    \item \textbf{Risk sharing:} Allows savers to spread investments across various assets, reducing risk exposure.
    \item \textbf{Liquidity:} Enables conversion of assets into cash easily, encouraging participation in investment markets.
    \item \textbf{Information:} Prices of financial assets reflect information about firms and projects, directing funds to their most productive uses.
\end{itemize}
These functions facilitate efficient allocation of savings into investments, supporting capital accumulation and growth.



\textbf{4. Crowding Out Effect}

\textbf{Question:}  
Explain how a government budget deficit affects the market for loanable funds, the equilibrium interest rate, and private investment.

\textbf{Solution:}  
\begin{itemize}
    \item A government deficit increases demand for loanable funds (shifts demand rightward).
    \item As a result, the \textbf{real interest rate rises}.
    \item Higher interest rates discourage private firms from borrowing to invest, leading to \textbf{crowding out}.
\end{itemize}

\textbf{Explanation:}
\begin{enumerate}
    \item In equilibrium, savings = investment (loanable funds market).
    \item When government borrows more, total demand for funds increases.
    \item Supply of funds is relatively fixed in the short run.
    \item The increased demand drives up interest rates, which makes private borrowing more expensive.
\end{enumerate}



\textbf{5. Phases of the Business Cycle}

\textbf{Question:}  
List the four phases of the business cycle and describe the typical movements of GDP, unemployment, and inflation during each phase.

\textbf{Solution:}  
\begin{enumerate}
    \item \textbf{Expansion:} Output, income, and employment rise; unemployment falls; inflation begins to rise.
    \item \textbf{Peak:} Output reaches its temporary maximum; inflation is often at its highest.
    \item \textbf{Recession:} Output and employment fall; unemployment rises; inflation declines.
    \item \textbf{Trough:} Economic activity bottoms out; inflation is low; recovery begins thereafter.
\end{enumerate}


\newpage
\section*{Part B – Quantitative Problems with Step-by-Step Solutions}

\textbf{6. Growth Rate Calculations}

\textbf{Question:}  
Real GDP in Canada increased from \$1.8 trillion in 2015 to \$2.2 trillion in 2020.  
(a) Compute the average annual growth rate.  
(b) Using the Rule of 70, estimate how long it would take for GDP to double at this rate.

\textbf{Solution:}

\underline{Step 1: Calculate growth rate}
\[
g = \left(\frac{2.2}{1.8}\right)^{1/5} - 1
\]
\[
g = (1.2222)^{0.2} - 1 = 0.041 = 4.1\%
\]

\underline{Step 2: Rule of 70}
\[
\text{Doubling time} = \frac{70}{g} = \frac{70}{4.1} \approx 17.1 \text{ years}
\]

\textbf{Interpretation:}  
At 4.1\% annual growth, real GDP would double in about 17 years.



\textbf{7. Per Capita Growth Rate}

\textbf{Question:}  
Country A’s real GDP grows by 5\% per year, while its population grows by 2\% per year.  
(a) Find the growth rate of real GDP per capita.  
(b) Interpret your result.

\textbf{Solution:}
\[
\text{Real GDP per capita growth} = 5\% - 2\% = 3\%
\]
\textbf{Interpretation:}  
On average, citizens’ living standards (output per person) rise by 3\% annually.

\newpage

\textbf{8. Market for Loanable Funds}

\textbf{Question:}  
Suppose demand for loanable funds is \( D = 800 - 50r \) and supply is \( S = 200 + 100r \), where \(r\) is the real interest rate (in \%).  
(a) Find equilibrium \(r\) and \(Q\).  
(b) If government borrowing increases by 100 (shifts demand right), find the new equilibrium.

\textbf{Solution:}

\underline{Step 1: Equilibrium condition}
\[
D = S \Rightarrow 800 - 50r = 200 + 100r
\]
\[
600 = 150r \Rightarrow r = 4\%
\]
\[
Q = 200 + 100(4) = 600
\]

\textbf{Equilibrium:} \(r^* = 4\%, Q^* = 600\)

\underline{Step 2: With government borrowing (+100 demand)}
\[
D' = 900 - 50r
\]
\[
900 - 50r = 200 + 100r \Rightarrow 700 = 150r \Rightarrow r' = 4.67\%
\]
\[
Q' = 200 + 100(4.67) = 667
\]

\textbf{Interpretation:}  
Higher government borrowing raises interest rates and reduces private investment due to crowding out.

\begin{center}
\rule{0.5\textwidth}{0.4pt}\\
\textit{End of Practice Problem Set}
\end{center}

\end{document}
